\section{Memory Windowing Experiment}
\label{appendix:infra_experiment_details}

\subsection{Experimental Setup}

The experiment compares two execution schedules for one standalone MoS memory block.
Whole-pack execution produces expert projections for all tokens at once; windowed execution groups complete steps, performs their projections, state updates, and readout in sequence, and recomputes internal intermediates during backward.
Both use the same equations, weights, inputs, precision, and update kernel.
Measurement covers both routers, all expert projections, memory writes and reads, RMSNorm, SoftCapSiLU, fusion, and backward through these computations.

\begin{table}[htbp]
    \centering
    \caption{Configuration of the MoS memory-block experiment.}
    \label{tab:infra_window_configuration}
    \small
    \renewcommand{\arraystretch}{1.12}
    \begin{tabular}{ll}
        \toprule
        Configuration & Value \\
        \midrule
        Hidden dimension & 2048 \\
        Memory experts & 4 \\
        Memory heads per expert / head dimension & 16 / 256 \\
        Router projection rank & 128 \\
        Decay-gate / output-gate projection rank & 256 / 256 \\
        Parameter and projection / memory-state precision & BF16 / FP32 \\
        Writer chunk & 64 tokens \\
        Step length & 2048 tokens \\
        Window budget & \makecell[l]{2048 / 4096 tokens,\\or 1 / 2 complete steps} \\
        \bottomrule
    \end{tabular}
\end{table}

Each pack contains two equal-length trajectories: 8K, 16K, and 32K packs comprise two 4K, 8K, and 16K trajectories, respectively.
Each trajectory starts from an independent initial state.
All measured steps commit state, while reads within a step use its incoming snapshot.
Windows do not cross trajectory boundaries.
Weights are randomly initialized and hidden-state inputs are synthetic, with fixed random seed 20261015; input hidden states have standard deviation 1.0, while initial states and backward cotangents have standard deviation 0.01.
This workload measures the computation and activation costs of the memory branch.

Memory updates use a fixed chunked-WY implementation.

\subsection{Measurement Method}

Each pack length and window budget define a paired experiment.
Both variants receive 5 warmups and 20 measured trials; successive pairs alternate between whole-pack$\rightarrow$windowed and windowed$\rightarrow$whole-pack order, yielding 240 measured samples.
Each window budget has an independently measured whole-pack baseline, accounting for the small baseline-time differences between rows at the same length.

CUDA events measure forward-plus-backward (F+B) time, including block execution, inner products of the output and final state with fixed cotangents, and backward through the scalar objective.
Input preparation, schedule construction, and compilation occur outside the timed region.
Peak memory is PyTorch peak allocated, including block parameters, inputs and cotangents, outputs, gradients, and both prepared schedules.
This metric measures live tensor allocations, distinct from memory reserved by the caching allocator.
The main table reports medians; percentage changes are computed from unrounded medians.

\subsection{Full Results}

\begin{table}[htbp]
    \centering
    \caption{Full memory and time results for the MoS memory block.
    Each row is an independently measured whole-pack/window pair; values are medians.}
    \label{tab:infra_window_full_results}
    \small
    \setlength{\tabcolsep}{5pt}
    \renewcommand{\arraystretch}{1.12}
    \begin{tabular}{rrrrrr}
        \toprule
        \makecell{Pack length\\(tokens)} & \makecell{Window budget\\(tokens)} & \makecell{Whole-pack\\memory (GiB)} & \makecell{Windowed\\memory (GiB)} & \makecell{Whole-pack\\F+B (ms)} & \makecell{Windowed\\F+B (ms)} \\
        \midrule
        8192 & 2048 & 12.37 & 4.85 & 102.1 & 122.9 \\
        8192 & 4096 & 12.37 & 8.36 & 102.1 & 117.4 \\
        16384 & 2048 & 24.19 & 5.23 & 211.6 & 246.1 \\
        16384 & 4096 & 24.19 & 8.77 & 211.6 & 235.1 \\
        32768 & 2048 & 47.83 & 6.03 & 462.7 & 493.5 \\
        32768 & 4096 & 47.83 & 9.49 & 463.0 & 470.6 \\
        \bottomrule
    \end{tabular}
\end{table}

As pack length increases from 8K to 32K, whole-pack peak memory grows from 12.37 GiB to 47.83 GiB, while 4096-token windows grow from 8.36 GiB to 9.49 GiB.
Windowing bounds the lifetime of projections and intermediates, while retained outputs and boundary states still grow with length.
Across the three lengths, the median time increases for 4096-token windows are 15.0\%, 11.1\%, and 1.6\%, respectively, showing that windowing overhead depends on the workload.

\subsection{Measurement Variability}

Table~\ref{tab:infra_window_variability} reports medians, interquartile intervals [Q25, Q75], and sample standard deviations, all in ms.

\begin{table}[htbp]
    \centering
    \caption{F+B timing variability in milliseconds.
    Intervals are the 25th and 75th percentiles; standard deviations use all measured samples.}
    \label{tab:infra_window_variability}
    \small
    \setlength{\tabcolsep}{4pt}
    \renewcommand{\arraystretch}{1.12}
    \begin{tabular}{rrlll}
        \toprule
        \makecell{Pack length\\(tokens)} & \makecell{Window budget\\(tokens)} & \makecell[l]{Whole-pack median\\{[Q25, Q75]}} & \makecell[l]{Windowed median\\{[Q25, Q75]}} & \makecell[l]{Standard deviation\\(whole / windowed)} \\
        \midrule
        8192 & 2048 & 102.10 [102.08, 102.12] & 122.87 [122.86, 122.88] & 0.03 / 0.02 \\
        8192 & 4096 & 102.11 [102.10, 102.12] & 117.42 [117.41, 117.45] & 0.05 / 0.03 \\
        16384 & 2048 & 211.55 [211.43, 212.02] & 246.11 [246.02, 246.42] & 37.19 / 0.26 \\
        16384 & 4096 & 211.61 [211.45, 211.87] & 235.07 [234.98, 235.39] & 0.30 / 0.27 \\
        32768 & 2048 & 462.71 [462.37, 463.87] & 493.48 [492.91, 494.38] & 22.22 / 0.94 \\
        32768 & 4096 & 462.96 [462.42, 463.68] & 470.58 [469.84, 471.46] & 0.76 / 38.01 \\
        \bottomrule
    \end{tabular}
\end{table}

All samples are retained, including whole-pack trials at 16K/2048 (378.03 ms) and 32K/2048 (562.40 ms), and a windowed trial at 32K/4096 (640.58 ms).
The 1.6\% reported in the chapter therefore describes the median time difference; individual trials exhibit variability.
